\section{The labelled tournament and nested predicates}
\label{sec:construction}

We construct the recursive predicates and the blocks that will witness their
block sensitivity at zero. The sensitivity estimates require control at
every input, so we also define the ordinary and joint profiles passed from
one level to the next.

Fix integers \(d\ge1\) and \(M\ge3\), and put
\[
 L=d+2,\qquad k=2M^2+1,\qquad r=\lceil\sqrt M\rceil,
 \qquad t=16.
\]
We will choose \(M\) as a function of \(d\) in
Section~\ref{sec:separation}.
A tournament is an orientation of every pair of distinct vertices:
exactly one of \(i\to j\) and \(j\to i\) holds.
Use a tournament on \([k]\) in which every vertex has outdegree \(M^2\).
For example, place the vertices cyclically and direct an edge from
each vertex to the next \(M^2\) vertices.

The tournament rows and selected blocking coordinates adapt the architecture
of Meiburg~\cite[Sections~3.1 and~5.2--5.5]{Meiburg2026}.
The specific label property needed for the ordinary-sensitivity argument is
the following elementary consequence of independent random labels.

\begin{lemma}\label{lem:labelling}
The edges can be labelled by numbers \(a(i,j)\in[r]\) so that there is
no set \(I\subseteq[k]\) of size \(t\) admitting choices
\(m_j\in[r]\), \(j\in I\), with
\[
 a(i,j)=m_j \qquad\text{for every edge }i\to j\text{ within }I.
\]
\end{lemma}
\begin{proof}
Choose all edge labels independently and uniformly from \([r]\).
For fixed \(I\) and fixed choices \((m_j)_{j\in I}\), the probability
of all the stated equalities is \(r^{-\binom t2}\).
There are at most \(k^t r^t\) choices to consider, so the probability
of any violation is at most
\[
 k^{16}r^{-104}\le (3M^2)^{16}(\sqrt M)^{-104}
 =3^{16}M^{-20}<1.
\]
Here \(k\le3M^2\), \(r\ge\sqrt M\), and \(M\ge3\).
Thus a labelling with the required property exists.
\end{proof}

Fix such a labelling for the remainder of the construction.
For \(0\le\ell\le d\), set
\[
 H_\ell=d+1-\ell,\qquad N_\ell=LM(kr)^\ell.
\]
We construct predicates \(P_{\ell,q}:\{0,1\}^{N_\ell}\to\{0,1\}\),
\(1\le q\le H_\ell\), that are nested in the sense that
\[
 P_{\ell,1}(x)\ge P_{\ell,2}(x)\ge\cdots
 \ge P_{\ell,H_\ell}(x)\quad\text{for every }x.
\]
Throughout, nesting concerns the predicate index, not the input bits.

\subsection{Initial predicates}
Partition \([LM]\) into \(M\) blocks \(E_1,\ldots,E_M\), each of
size \(L\). Let \(z_j\) be the string that is one precisely on \(E_j\).
Write \(\operatorname{dist}(x,z)\) for Hamming distance and set
\[
 D(x)=\min_{j\in[M]}\operatorname{dist}(x,z_j),\qquad
 P_{0,q}(x)=\mathbf1\{D(x)\le H_0-q\}.
\]
These predicates are nested. Their radii are the integers from \(d\)
down to zero, while distinct centres have distance \(2L=2d+4\).

\subsection{Recursive clauses}
Suppose \(1\le\ell\le d\). Partition the input into \(kr\) disjoint
children \(x_{j,c}\in\{0,1\}^{N_{\ell-1}}\), arranged in \(k\) rows
\(j\in[k]\), each with \(r\) positions \(c\in[r]\).
For this level put \(H=H_\ell\) and \(Q=H+1=H_{\ell-1}\).
For \(q\in[H]\), let clause \(i\) at index \(q\) assert
\begin{equation}\label{eq:predicate-definition}
 \begin{aligned}
 P_{\ell-1,Q}(x_{i,c})&=1 &&\text{for all }c\in[r],\\
 P_{\ell-1,Q-q}(x_{j,a(i,j)})&=0 &&\text{for all }i\to j.
 \end{aligned}
\end{equation}
Call the first line its \emph{target conditions} and the second line
its \emph{gate conditions}. Figure~\ref{fig:clause} depicts the target row
and one outgoing gate. Define \(P_{\ell,q}(x)=1\) if at least
one clause is satisfied, and zero otherwise. All indices are valid:
\(1\le Q-q\le H<Q\).
Within a single clause, all \(r+M^2\) conditions use distinct children.
Indeed, targets use row \(i\), each gate uses another row, and there
is only one edge from \(i\) to any given row. Consequently a change
of one raw input bit can affect at most one condition of that clause.
The same bit can affect several clauses, which will cause no problem
when taking upper bounds by unions of sets of bits.

By induction, increasing \(q\) makes each gate condition stronger:
the index \(Q-q\) decreases, and requiring a larger nested predicate
to be zero is more restrictive. Targets do not change. Thus each
clause, and hence each \(P_{\ell,q}\), is nested as asserted.
This recursion defines total Boolean functions on exactly
\(N_\ell=krN_{\ell-1}\) bits.

\begin{figure}[htbp]
\centering
\begin{tikzpicture}[>=Latex, every node/.style={font=\small},
  condition/.style={draw,rounded corners=2pt,inner sep=8pt,align=center}]
 \node[condition] (target) at (0,0)
   {Row \(i\): all \(r\) children\\[2pt]
    \(P_{\ell-1,Q}(x_{i,c})=1\)};
 \node[condition] (gate) at (6,0)
   {Row \(j\): child \(a(i,j)\)\\[2pt]
    \(P_{\ell-1,Q-q}(x_{j,a(i,j)})=0\)};
 \draw[->] (target) -- node[above]{\(i\to j\)} (gate);
 \node[below=5pt of target,align=center] {One target condition\\per child in row \(i\)};
 \node[below=5pt of gate,align=center] {One gate condition\\per outgoing neighbour \(j\)};
\end{tikzpicture}
\caption{One recursive clause. The edge label selects a single child in
each outgoing neighbour's row. Raising \(q\) strengthens the gates;
all targets use the strongest child predicate \(Q\).}
\label{fig:clause}
\end{figure}

\subsection{Sensitivity profiles and the initial bounds}
For \(b\in\{0,1\}\), define
\[
 S_{\ell,b}=\max_{\substack{1\le q\le H_\ell,\ x\in\{0,1\}^{N_\ell}\\
                           P_{\ell,q}(x)=b}}
 \bigl|\{a\in[N_\ell]:P_{\ell,q}(x^a)=1-b\}\bigr|.
\]
We also keep track of simultaneous changes of two distinct predicates:
\[
 J_{\ell,b}=\max_{\substack{1\le q<q'\le H_\ell,\ x\in\{0,1\}^{N_\ell}\\
                         P_{\ell,q}(x)=P_{\ell,q'}(x)=b}}
 \bigl|\{a\in[N_\ell]:P_{\ell,q}(x^a)=P_{\ell,q'}(x^a)=1-b\}\bigr|.
\]
A maximum over no cases is defined to be zero; in particular
\(J_{d,0}=J_{d,1}=0\), since \(H_d=1\).
The two predicates in the definition of \(J_{\ell,b}\) are evaluated
on the same input and changed by the same raw input bit.

\begin{lemma}\label{lem:base}
The initial sensitivity profiles satisfy
\begin{equation}\label{eq:base-profile}
 S_{0,0}\le L,\qquad S_{0,1}\le LM,\qquad
 J_{0,0}=J_{0,1}=0.
\end{equation}
\end{lemma}
\begin{proof}
Fix a predicate of radius \(R=H_0-q\in\{0,\ldots,d\}\) and an
input \(x\) at which it is zero. If flipping one bit makes it one,
some centre must be at distance exactly \(R+1\) from \(x\).
There is at most one centre within distance \(d+1\) of \(x\):
two such centres would have distance at most \(2d+2<2L\).
Thus every successful bit flip must move towards the same centre,
giving at most \(R+1\le d+1\le L\) possibilities.
The bound on \(S_{0,1}\) is simply the number \(LM\) of input bits.

For adjacent inputs \(x,y\), the triangle inequality gives
\(D(x)\le D(y)+1\), and the reverse inequality follows by symmetry.
Hence \(|D(x)-D(y)|\le1\). To change both predicates with distinct
integer radii from zero to one, or both from one to zero, would require
\(D\) to change by at least two. Both joint sensitivities are therefore
zero.
\end{proof}

\subsection{Disjoint blocks at zero}
\begin{lemma}\label{lem:blocks}
At every level \(0\le\ell\le d\), all predicates are zero at the
all-zero input. There are \(Mk^\ell\) pairwise disjoint nonempty blocks
of coordinates, each of size \(Lr^\ell\), such that flipping any one
block at zero makes every predicate at that level equal to one.
\end{lemma}
\begin{proof}
At level zero, \(D(0)=L>d\), and the blocks \(E_j\) have the stated
property because \(0^{E_j}=z_j\).
Suppose the claim holds at level \(\ell-1\). At the all-zero input,
every target predicate is zero, so every clause fails at level \(\ell\).

Enumerate the preceding-level block family as
\(B_1,\ldots,B_{Mk^{\ell-1}}\).
For each row \(i\) and each family index \(h\), take the union of the
copies of \(B_h\) in all \(r\) children of row \(i\).
Flipping this union makes the strongest predicate in each of those
children equal to one, so all targets for clause \(i\) hold.
All gate children of clause \(i\) lie in other rows and remain zero,
so its gates hold for every \(q\).
Thus every \(P_{\ell,q}\) is one after the flip.
Within a row these unions are disjoint because the preceding-level
blocks are disjoint in every child; unions from different rows are
also disjoint. There are \(k\cdot Mk^{\ell-1}=Mk^\ell\) of them,
each of size \(r\cdot Lr^{\ell-1}=Lr^\ell\), as required.
\end{proof}
